\section{PairSelect}
\label{sec:method}
\begin{table}[!b]
\centering\small
\caption{Actual token counts entering blocks 7--12. Headers specify nominal deletion; CLS is included when present. Video encoders count tokens per clip; CLIP counts per image.}
\label{tab:budgets}
\setlength{\tabcolsep}{3pt}
\begin{tabular}{@{}lrrrr@{}}
\toprule
Visual encoder + detector & Dense & 20\% & 40\% & 60\%\\
\midrule
\tableinput{tables/theory_rows}
\bottomrule
\end{tabular}
\end{table}
\begin{table*}[!t]
\centering\small
\caption{Four-system quality at every budget (0--100). Each cell reports the indicated two frame-level metrics; bold identifies the primary metric, not a ranking. On XD, the video systems use trapezoidal PR-AUC and DSANet uses step AP. CI is the paired 95\% interval for the primary-metric change at 40\% deletion, in percentage points.}
\label{tab:quality}
\setlength{\tabcolsep}{4pt}
\begin{tabular}{@{}lrrrrc@{}}
\toprule
System & Dense & Delete 20\% & Delete 40\% & Delete 60\% & CI of $\Delta_{40}$\\
\midrule
\tableinput{tables/quality_rows}
\bottomrule
\end{tabular}
\end{table*}


\subsection{Direct insertion into frozen systems}
Let $E_{1:b}$ and $E_{b+1:L}$ be the prefix and suffix of a pretrained visual encoder, $R$ its readout, and $H$ its temporal detection head. For scheduled visual inputs $\{v_t\}_{t=1}^{M}$,
\begin{equation}
 z_t^{(r)}=R\!\left(E_{b+1:L}\!\left(S_r(E_{1:b}(v_t))\right)\right),
 \quad \hat{\boldsymbol y}^{(r)}=H(\{z_t^{(r)}\}_{t=1}^{M}).
 \label{eq:pipeline}
\end{equation}
Only the selection operator $S_r$ is inserted, with retention ratio $r$. All four visual encoders use $L=12$, $b=6$. The original input schedule, feature dimension, temporal aggregation, and head weights are preserved. For the three video backbones, each UR-DMU head is trained once on the corresponding dense training features; DSANet uses the authors' released detector weights. \emph{Training-free} refers to inserting and changing the compression operator, not to training the original systems.

Local redundancy motivates comparisons among neighboring patches, while selection avoids introducing additional vector mixtures at the insertion point. Activation magnitude provides an inexpensive preference, inspired by magnitude-aware sampling~\cite{fayyaz2022ats}. It is not an anomaly probability~\cite{darcet2024registers}; the local quota prevents an unconstrained global ranking from concentrating the entire budget in a few regions.

\subsection{Local quotas and member priority}
For VideoMAE, VideoMAEv2, and CLIP, spatially adjacent horizontal patches form pairs ordered by native coordinates. Each consecutive group contains at most ten patch tokens (five pairs). A group of $n_g$ tokens receives
\begin{equation}
 q_g(r)=\lfloor rn_g+\tfrac12\rfloor,
 \qquad r\in\{0.8,0.6,0.4\}.
 \label{eq:quota}
\end{equation}
For pair $j$ with members $(x_{j,0},x_{j,1})$, define
\begin{equation}
 a_j=\arg\max_{u\in\{0,1\}}\norm{x_{j,u}},\quad
 s_j=\max_{u\in\{0,1\}}\norm{x_{j,u}}.
 \label{eq:score}
\end{equation}
Let $w_j=x_{j,a_j}$ and $\ell_j=x_{j,1-a_j}$. If $\pi$ sorts the group's $m_g$ pairs by descending $s_j$, the priority sequence is
\begin{equation}
 \mathcal P_g=(w_{\pi_1},\ldots,w_{\pi_{m_g}},
                 \ell_{\pi_1},\ldots,\ell_{\pi_{m_g}}).
 \label{eq:priority}
\end{equation}
The first $q_g$ entries are retained and restored to native index order before gathering. Ties preserve native order. At 20\% and 40\% deletion, every pair in a full group retains a representative, with three or one additional members; 60\% deletion retains four of five representatives. Selected vectors are neither averaged nor rescaled. CLIP applies the rule independently per image and always retains CLS, yielding the counts in Table~\ref{tab:budgets}.

\subsection{Structure-aware adaptation}
For TimeSformer, the selection unit is a complete spatial trajectory of $T$ tokens, scored by its maximum member norm. Each group of ten trajectories retains its top $q_g$. The current bridge rounds the total spatial count down to a multiple of native width $W$, removing excess quota from tail groups:
\begin{equation}
 P'=W\left\lfloor\frac{\sum_g q_g(r)}{W}\right\rfloor,
 \qquad K=TP'+1.
 \label{eq:trajectory}
\end{equation}
With $P=196$, $T=8$, and $W=14$, the three budgets retain 154, 112, and 70 trajectories plus CLS. Actual token deletion is 21.41\%, 42.83\%, and 64.24\%; this alignment constraint is implementation-specific. Video-encoder bridges use a fixed quota-slot skeleton with approximate spatial anchors while retaining trajectory time indices. CLIP gathers the already position-embedded patch states in native order and preserves its CLS readout. These adaptations keep the detector interface unchanged.
